% 原书分章 PDF 第 9—11 页：第 11.8—11.10 节与图11.1。
\section{TMLE 的统计性质}

累积发生率的 TMLE 估计量具有双重稳健性（double robustness）：只要 $Q_n$ 一致估计 $Q_0$，或 $G_n$ 一致估计 $G_0$，该估计量就是一致的。在基于 TMLE 的估计量所要求的常规经验过程条件和收敛速率条件下，TMLE 估计量具有渐近线性。如果处理机制和删失机制完全已知，$\psi_n^*$ 的影响函数为 $D^*(Q^*,G_0)$，其中 $Q^*$ 是 $Q_n^*$ 的极限，可能存在模型设定错误。在这种情形下，可以用
\[
 \sigma_n^2=\sum_{i=1}^nD^*(Q_n^*,G_0)(O_i)^2
\]
一致估计 $n^{1/2}(\psi_n^*-\psi_0)$ 的渐近方差。然而，在实践中，处理机制与删失机制通常未知；如果同时 $Q_n^*$ 对 $Q_0$ 不一致，渐近方差就更加复杂。即便如此，若在这种情形中，$G_n$ 是参数模型 $\mathcal G$ 内的一个渐近有效估计量，仍可以将 $\sigma_n^2$ 用作渐近方差的保守估计。

可以为 $\psi_n^*$ 构造以下置信区间：
\[
 \bigl(
 \psi_n^*-z_{1-\alpha/2}\sigma_n n^{-1/2},\,
 \psi_n^*+z_{1-\alpha/2}\sigma_n n^{-1/2}
 \bigr),
\]
其中 $z_\beta$ 是标准正态分布的 $\beta$ 分位数。类似地，给定固定值 $\psi^\circ\in(0,1)$，可以构造原假设 $\psi_0=\psi^\circ$ 的双侧检验，其渐近检验水平不超过 $\alpha$：当
\[
 \left|n^{1/2}(\psi_n^*-\psi^\circ)/\sigma_n\right|>z_{1-\alpha/2}
\]
时，拒绝原假设。基于影响函数的方差估计有一个很吸引人的特点：若某个量是多个估计量的函数，可以方便地得到它的渐近方差形式。在筛选分析中，这一点很有用，因为由此可以方便地构造基因型特异疫苗效力和疫苗筛选效应的置信区间及假设检验。

\section{HVTN 505 HIV 疫苗筛选分析}

我们分析了 HVTN 505 研究的数据。这是一项近期开展的预防性 HIV 疫苗效力 II 期试验，参与者按 $1:1$ 随机分配，接受候选疫苗或安慰剂（Hammer 等，2013）。研究在招募时，还收集了参与者风险行为的其他信息，包括：性别（男性、女性、跨性别男性、跨性别女性）、种族（白人、黑人、其他）、身体质量指数（BMI）、药物使用情况（大麻、可卡因、poppers、speed、MDMA、其他娱乐性药物）、饮酒情况（从不、每周不足一次、每周 1—2 天、每周 3—6 天、每天）、性传播疾病状况（梅毒、疱疹、生殖器溃疡、淋病），以及性行为风险（自报性伴侣人数、自报与男性发生无保护插入方肛交、自报与男性发生无保护接受方性行为、此前构建的行为风险评分）。接受最后一剂疫苗后，参与者每 3 个月参加一次访视，接受 HIV 检测。我们重点分析修正意向性治疗队列（modified intent-to-treat cohort），其中包括 2504 名参与者。有关试验设计的更多信息，可参阅原始发表论文。

我们研究了该疫苗在 HIV 包膜蛋白 V2 环的第 169 号氨基酸位点，是否表现出筛选效应。选择这一位点，是因为一款与 HVTN 505 疫苗设计相似的疫苗，曾在该位点表现出筛选效应（Rolland 等，2012）。在 HVTN 505 中，共发生 17 例 169 位点匹配感染和 30 例 169 位点不匹配感染。我们使用 TMLE，估计每个预定访视时间窗内，疫苗组与安慰剂组中 169 位点匹配感染和不匹配感染的累积发生率。再利用这些度量，估计疫苗对每类感染的效力，以及该位点的疫苗筛选效应。

我们使用超级学习器，估计迭代条件均值和条件删失分布。两者使用相同的算法库，共包含 13 个算法：一个未经调整的均值；三个仅含主项的广义线性模型，分别仅使用行为风险评分、仅使用性行为风险变量、仅使用与结局相关性最高的五个变量；三个仅含主项、基于 AIC 的逐步回归，分别使用全部变量、仅使用性行为风险变量、仅使用相关性最高的五个变量；三个随机森林，分别使用全部变量、仅使用药物或酒精使用变量、仅使用性行为风险变量；以及三个梯度提升机，分别使用全部变量、仅使用药物或酒精使用变量、仅使用全部性行为风险变量。

\begin{figure}[htbp]
 \centering
 % 图11.1：原始四面板矢量数据、置信区间及截断箭头均原样保留；只替换英文标签。
\begin{tikzpicture}[x=.45mm,y=-.45mm]
 \node[anchor=north west,inner sep=0pt] at (0,0)
  {\includegraphics[width=150.3mm,height=96.3mm]{figures/ch11_sieve_data.pdf}};
 \node[font=\small,rotate=90] at (5.9503,45.9548) {累积发生率};
 \node[anchor=west,font=\fontsize{7.2}{7.7}\selectfont] at (35.9112,17.3084) {安慰剂};
 \node[anchor=west,font=\fontsize{7.2}{7.7}\selectfont] at (35.9112,24.2791) {疫苗};
 \node[font=\small] at (87.8391,6.5254) {氨基酸 169 位点匹配};
 \node[font=\small,rotate=90] at (173.3625,45.9543) {累积发生率};
 \node[anchor=west,font=\fontsize{7.2}{7.7}\selectfont] at (203.9068,17.3085) {安慰剂};
 \node[anchor=west,font=\fontsize{7.2}{7.7}\selectfont] at (203.9068,24.2791) {疫苗};
 \node[font=\small] at (255.8328,6.5254) {氨基酸 169 位点不匹配};
 \node[font=\small] at (87.8429,208.2375) {末次疫苗接种后的月数};
 \node[font=\small,rotate=90] at (5.9508,147.5245) {疫苗效力};
 \node[anchor=west,font=\fontsize{7.2}{7.7}\selectfont] at (40.5308,114.8452) {169 位点匹配};
 \node[anchor=west,font=\fontsize{7.2}{7.7}\selectfont] at (40.5308,121.8159) {169 位点不匹配};
 \node[font=\small] at (255.8344,208.2375) {末次疫苗接种后的月数};
 \node[font=\small,rotate=90] at (173.3615,147.5202) {筛选效应};
\end{tikzpicture}

 \caption{HVTN-505 的 TMLE 分析结果。}
 \label{fig:ch11-sieve}
\end{figure}

分析结果见图 11.1。上方一行展示了两种处理下的 HIV 感染累积发生率，其中，感染病毒在目标氨基酸位点上，与疫苗所编码的病毒株匹配（左图）或不匹配（右图）。总体上，不匹配感染的发生率略高于匹配感染；但在两种情形中，两个处理组的曲线都没有明显分离。左下图中的疫苗效力度量也反映了这一点：该疫苗似乎对两类感染都没有保护效力。右下面板给出筛选效应的估计，进一步表明，没有证据说明疫苗对匹配感染和不匹配感染具有不同的效力。

我们还使用 Aalen--Johansen 估计量进行了分析，以考察调整协变量是否会影响结果。在接受最后一剂疫苗 18 个月后，TMLE 估计的筛选效应为 0.83（95\% 置信区间：0.31—2.23），Aalen--Johansen 估计的筛选效应为 0.82（95\% 置信区间：0.23—2.91）。协变量调整没有影响筛选效应的点估计，但对该点估计的不确定性估计有明显影响。TMLE 筛选效应估计的估计方差，是 Aalen--Johansen 筛选效应估计的 60\%，因而 TMLE 估计对应的置信区间明显更窄。

\section{讨论}

本章说明了如何使用 TMLE，估计累积发生率、基因型特异的疫苗效力，以及疫苗筛选效应。

\begin{sourcebox}
这种靶向筛选分析方法具有若干优点。协变量调整允许删失依赖基线协变量，从而可以放宽独立删失假设。此外，协变量调整还能够提高效率。因此，与标准的、未经调整的筛选分析方法相比，我们可以得到偏差更小、方差也更小的疫苗效力和筛选效应估计量。
\end{sourcebox}

对 HVTN 505 数据的分析表明，这些潜在优势确实能够在真实数据应用中实现。调整 HVTN 505 试验所收集的大量风险行为变量，使筛选效应估计对应的置信区间明显变窄。疫苗筛选效应的靶向学习，看来是这一领域一个值得发展的方向。
